\documentclass[french]{book}
\usepackage{teach}
\usepackage{verbatim}
\usepackage{amsmath,amsfonts,amssymb}
\usepackage{pgfplots}
\usepackage{mwe}
\RequirePackage{graphicx}
\RequirePackage{ifpdf}
 \ifpdf
   \DeclareGraphicsRule{*}{mps}{*}{}
 \fi


  \usepackage{fancyhdr}

  \usepackage{amsmath}
  \usepackage{amssymb}
  \usepackage{amsfonts}
  \usepackage{amsthm}
  \usepackage{mathrsfs}

  \usepackage{array}
  \usepackage{multicol}
  \setlength{\multicolsep}{3pt}

  \usepackage{graphicx}
  \graphicspath{{Figures/}}
  \usepackage{pstricks,pst-plot,pst-text,pst-tree,pst-eucl}
  \usepackage[all]{xy}

  \usepackage{fancybox}
  \usepackage{color}

%  \usepackage[frenchb]{babel}
  \usepackage{hyperref}

\usepackage{subfig}
\pgfplotsset{compat=newest}
\usepackage[all]{xy}
%\usepackage{geometry}
%\geometry{top=1cm, bottom=1cm, left=2cm, right=2cm}
\usepackage[utf8]{inputenc}
\usepackage[french]{babel}
\redefinecolor{thm}{title}{bgcolor}{alizarin}
\redefinecolor{prop}{title}{bgcolor}{meatbrown}
\redefinecolor{defn}{title}{bgcolor}{olivine}
\definecolor{alizarin}{rgb}{0.82, 0.1, 0.26}
\definecolor{meatbrown}{rgb}{0.9, 0.72, 0.23}
\definecolor{olivine}{rgb}{0.6, 0.73, 0.45}
\usepackage{mathrsfs}
%%
\newcommand{\R}{\mathbb {R}}
%\newcommand{\C}{\mathds {C}}
\newcommand{\N}{\mathbb {N}}
\newcommand{\Z}{\mathbb {Z}}
\newcommand{\Q}{\mathbb {Q}}
\newcommand{\D}{\mathbb {D}}
%%
\newcommand{\se}{\geq}
\newcommand{\ie}{\leq}
\newcommand{\tm}{\times}
%\newcommand{\ell}{l}
%\newcommand{\ell'}{l'}
%%
\newcommand{\limn}{\lim_{n\rightarrow +\infty}}
\newcommand{\limo}[1][x]{\ds\lim_{#1\rightarrow 0}}
\newcommand{\limite}[2][x]{\ds\lim_{#1\rightarrow #2}}
\newcommand{\limpinf}[1][x]{\ds\lim_{#1\rightarrow +\infty}}
\newcommand{\limminf}[1][x]{\ds\lim_{#1\rightarrow -\infty}}
\newcommand{\limiteg}[2][x]{\ds\lim_{#1\xrightarrow{<}#2}}
\newcommand{\limited}[2][x]{\ds\lim_{#1\xrightarrow{>}#2}}
%%
\newcommand{\vect}[1]{\overrightarrow{#1}}
\newcommand{\oij}{(\textit{O},{\overrightarrow{i}},{\overrightarrow{j}})}
%
\newcommand{\cp}[2]{%
        \begin{pmatrix}
        #1\\
        #2
        \end{pmatrix}%
        }
%
\newcommand{\calb}{\mathscr{B}}
\newcommand{\calc}{\mathscr{C}}
\newcommand{\cald}{\mathscr{D}}
\newcommand{\cale}{\mathscr{E}}
\newcommand{\calf}{\mathscr{F}}
\newcommand{\calp}{\mathscr{P}}
\newcommand{\cals}{\mathscr{S}}
\newcommand{\calt}{\mathscr{T}}
%
\newcommand{\suite}[1][u]{\left( #1_{n} \right)_{n\in\N}}
\newcommand{\suitep}[1][u]{\left( #1_{n} \right)_{n\in\N^{*}}}
\usetikzlibrary{arrows}

\newcommand{\poly}[1][a]{#1_0+#1_1x+\cdots+#1_nx^n}


\begin{document}
\tableofcontents
\chapter{Polynômes du second degré}
\section{Généralités sur les polynômes}
\subsection{Définition}
\begin{defn}
On appelle fonction polynôme toute fonction $f$ définie sur $\R$ par :
\[f(x)=\sum^n_{i=0} a_ix^i = a_0+a_1x+\cdots+a_nx^n \]
où $a_{0}$, $a_{1}$, $\ldots$, $a_{n}$ sont des réels donnés.
\end{defn}

\begin{exemple}
$x\mapsto -x^{3}-5x^{2}+7x-1$ est un polynôme. $x\mapsto \dfrac{x^{4}-4}{x^{2}+2}$ n'est pas un polynôme.
\end{exemple}%

\subsection{Propriétés (admises)}
\begin{prop}
\begin{enumerate}
\item Soit $P$ le polynôme défini par $P(x)=\poly{a}{n}$.

$P$ est le polynôme nul $\Longleftrightarrow$ $a_{0}=a_{1}=\cdots=a_{n}=0$.
\item Soient $P$ et $Q$ les polynômes définis par $P(x)=\poly{a}{n}$ et $Q(x)=\poly{b}{p}$ avec $a_{n}\neq0$ et $b_{p}\neq0$.
\[P=Q\Longleftrightarrow\left\lbrace
\begin{array}{l}
n=p \\
a_{0}=b_{0}\,;\, a_{1}=b_{1}\,;\,\cdots \,a_{n}=b_{n}
\end{array}
\right.\]
\end{enumerate}
\end{prop}
Conséquence : L'écriture d'un polynôme est unique.

\begin{exemple}
Si, pour tout $x\in\R$, $ax^3+bx^2+cx+d=2x^3-x+2$ alors $a=2$, $b=0$, $c=-1$ et $d=1$.
\end{exemple}
%
\subsection{Degré}
\begin{defn}
Soit $P$ un polynôme défini par $P(x)=\poly[a]$ avec $a_{n}\neq0$.

Le nombre $n$ est appelé degré de $P$.
\end{defn}

\begin{exemple}
$x\mapsto -x^{3}-5x^{2}+7x-1$ est un polynôme de degré 3.
$x\mapsto 3x-x^{5}$ est un polynôme de degré 5.
\end{exemple}%

\subsection{Racines d'une fonction}
\begin{defn}
Soit $f$ une fonction. On appelle racine de $f$ toute solution de l'équation $f(x)=0$.
\vspace{5mm}
\end{defn}

\begin{exemple}
1 est une racine de $x\mapsto -x^{3}-5x^{2}+7x-1$. 0 est une racine de $x\mapsto 3x-x^{5}$.
\end{exemple}
%%%%%%%%%%%%%%%%%%%%%%%%%%%%%%%%%%%%%%%%%%%%%%%%%%%%%%%%%%%%%%%%%%%%%%%%%%%%%%%%%%%%%
\section{Polynômes du second degré}
Dans tout le paragraphe, $P$ désigne un polynôme défini par $P(x)=ax^{2}+bx+c$ avec $a\neq0$.
\subsection{Forme canonique}
\begin{prop}[Propriété et définition]
Il existe des réels $\alpha$ et $\beta$ tels que :
\begin{center}
Pour tout $x\in\R$, $P(x)=a\left((x+\alpha)^{2}+\beta\right)$.
\end{center}
Cette écriture est appelée forme canonique de $P$.
\end{prop}

Remarque : On appelle parfois forme canonique l'écriture de $P$ sous la forme
\[P(x)=a(x+\alpha)^{2}+\gamma\]

\begin{demo}
\begin{align*}
P(x)&=ax^{2}+bx+c=a\left(x^{2}+\dfrac{b}{a}x+\dfrac{c}{a}\right)=a\left(\left(x+\dfrac{b}{2a}\right)^{2}-\dfrac{b^{2}}{4a^{2}}+\dfrac{c}{a}\right)\\
&=a\left(\left(x+\dfrac{b}{2a}\right)^{2}+\dfrac{-b^{2}+4ac}{4a^{2}}\right)
\end{align*}
\end{demo}

\begin{exemple}
Écrire la forme canonique du polynôme défini par $P(x)=2x^{2}+4x+6$.
\end{exemple}


$P(x)=2x^{2}+4x+6=2\left(x^{2}+2x+3\right)=2\left((x+1)^{2}-1+3\right)=2\left((x+1)^{2}+2\right)$


\subsection{Discriminant}
\begin{defn}
On appelle discriminant de $P$ le réel $\Delta$ défini par $\Delta=b^{2}-4ac$.
\vspace{5mm}
\end{defn}

\begin{exemple}
Calculer le discriminant du polynôme défini par $P(x)=2x^{2}+4x+6$.
\end{exemple}


$\Delta=4^{2}-4\times 2\times 6=16-48=-32$.


\subsection{Racines}
\begin{prop}
Les racines de $P$ peuvent être déterminée de la façon suivante :
\begin{itemize}
\item Si $\Delta<0$ alors $P$ n'a pas de racine réelle.
\item Si $\Delta=0$ alors $P$ admet une racine réelle : $-\dfrac{b}{2a}$.
\item Si $\Delta>0$ alors $P$ admet deux racines réelles : $\dfrac{-b-\sqrt{\Delta}}{2a}$ et $\dfrac{-b+\sqrt{\Delta}}{2a}$.
\end{itemize}
\end{prop}

\begin{demo}
\[P(x)=a\left(\left(x+\dfrac{b}{2a}\right)^{2}+\dfrac{-b^{2}+4ac}{4a^{2}}\right)=a\left(\left(x+\dfrac{b}{2a}\right)^{2}-\dfrac{\Delta}{4a^{2}}\right)\]
\begin{itemize}
\item Si $\Delta<0$ alors $\left(x+\dfrac{b}{2a}\right)^{2}-\dfrac{\Delta}{4a^{2}}>0$ donc pour tout $x$, $P(x)\neq0$.
\item Si $\Delta=0$ alors $P(x)=a\left(x+\dfrac{b}{2a}\right)^{2}$ donc $P(x)=0\Longleftrightarrow x=-\dfrac{b}{2a}$.
\item Si $\Delta>0$ alors $P(x)=a\left(\left(x+\dfrac{b}{2a}\right)-\dfrac{\sqrt{\Delta}}{2a}\right)\left(\left(x+\dfrac{b}{2a}\right)+\dfrac{\sqrt{\Delta}}{2a}\right)$

donc
\begin{align*}
P(x)=0&\Longleftrightarrow\left(x+\dfrac{b}{2a}\right)-\dfrac{\sqrt{\Delta}}{2a}=0 \quad \text{ou} \quad \left(x+\dfrac{b}{2a}\right)+\dfrac{\sqrt{\Delta}}{2a}=0\\
&\Longleftrightarrow x=\dfrac{-b+\sqrt{\Delta}}{2a} \quad \text{ou} \quad x=\dfrac{-b+\sqrt{\Delta}}{2a}
\end{align*}

\end{itemize}
\end{demo}

\begin{exemple}
Déterminer les racines de $P$ et $Q$ définis par
 \[P(x)=2x^{2}+4x+6 \quad \text{et} \quad Q(x)=-x^{2}-2x+1\]
\end{exemple}

Pour $P$ : $\Delta=-32<0$ donc $P$ n'a pas de racine.

Pour $Q$ : $\Delta=(-2)^{2}-4\times (-1)\times 1=4+4=8>0$ donc $Q$ admet deux racines :
\begin{align*}
x_{1}&=\dfrac{-(-2)-\sqrt{8}}{2\times (-1)}=\dfrac{2-2\sqrt{2}}{-2}=-1+\sqrt{2}\\
x_{2}&=\dfrac{-(-2)+\sqrt{8}}{2\times (-1)}=\dfrac{2+2\sqrt{2}}{-2}=-1-\sqrt{2}
\end{align*}



\subsection{Relation de Viète}
\begin{prop}
Si $P$ admet deux racines $x_{1}$ et $x_{2}$ alors 
$\left\lbrace 
\begin{array}[c]{l}
r_{1}+r_{2}=-\dfrac{b}{a} \\
r_{1}r_{2}=\dfrac{c}{a}
\end{array}
\right.$
\end{prop}

%\ifthenelse{\value{version}>0}{\newpage}{}
\begin{demo}
\begin{align*}
x_{1}+x_{2}&=\dfrac{-b-\sqrt{\Delta}}{2a}+\dfrac{-b+\sqrt{\Delta}}{2a}=\dfrac{-2b}{2a}=-\dfrac{b}{a}\\
x_{1}x_{2}&=\dfrac{-b-\sqrt{\Delta}}{2a}\times \dfrac{-b+\sqrt{\Delta}}{2a}=\dfrac{(-b)^{2}-\Delta}{4a^{2}}=\dfrac{b^{2}-b^{2}+4ac}{4a^{2}}=\dfrac{4ac}{4a^{2}}=\dfrac{c}{a}
\end{align*}
\end{demo}%

\subsection{Factorisation}
\begin{prop}
\begin{itemize}
\item Si $\Delta<0$ alors $P$ ne peut pas être factorisé.
\item Si $\Delta=0$ alors $P(x)=a\left(x+\dfrac{b}{2a}\right)^{2}$
\item Si $\Delta>0$ alors $P(x)=a(x-x_{1})(x-x_{2})$ où $x_{1}$ et $x_{2}$ sont les racines de $P$.
\end{itemize}
\end{prop}

%\ifthenelse{\value{version}=0}{\newpage}{}
\begin{demo}
Le premier résultat s'obtient en remarquant que si $P$ pouvait se factoriser, on aurait deux facteurs du premier degré auquel cas l'équation $P(x)=0$ admettrait au moins une solution, ce qui est contradictoire avec le résultat obtenu précédemment.

Les deux autres résultats ont été obtenus dans le cours de la démonstration précédente.
\end{demo}

\begin{exemple}
Factoriser $P$ et $Q$ définis par
 \[P(x)=2x^{2}+4x+6 \quad \text{et} \quad Q(x)=-x^{2}-2x+1\]
\end{exemple}


Pour $P$ : $\Delta=-32<0$ donc $P$ ne peut pas se factoriser.

Pour $Q$ : Les racines sont $-1+\sqrt{2}$ et $-1+\sqrt{2}$ donc $P(x)=-(x+1-\sqrt{2})(x+1+\sqrt{2})$.

%
\subsection{Signe}
\begin{prop}
\begin{itemize}
\item Si $\Delta<0$ alors
\begin{tabular}[c]{|c|ccc|} \hline
$x$ & $-\infty$ &  & $+\infty$ \\ \hline
Signe de $P(x)$ &  & Signe de $a$ & \\ \hline
\end{tabular}
\item Si $\Delta=0$ alors
\begin{tabular}[c]{|c|ccccc|} \hline
$x$ & $-\infty$ &  & $-\frac{b}{2a}$ &  & $+\infty$ \\ \hline
Signe de $P(x)$ &  & Signe de $a$ & $0$ & Signe de $a$ & \\ \hline
\end{tabular}
\item Si $\Delta>0$ alors

\vspace{1ex}
\begin{tabular}[c]{|c|ccccccc|} \hline
$x$ & $-\infty$ &  & $x_{1}$ &  & $x_{2}$ &  & $+\infty$ \\ \hline
Signe de $P(x)$ &  & Signe de $a$ & $0$ & Signe de $-a$ & $0$ & Signe de $a$ & \\ \hline
\end{tabular}

\vspace{1ex}
\hspace*{1cm}où $x_{1}$ et $x_{2}$ sont les racines de $P$ et $x_{1}<x_{2}$
\end{itemize}
\end{prop}

\begin{demo}
\begin{itemize}
\item Si $\Delta<0$ alors $\left(x+\dfrac{b}{2a}\right)^{2}-\dfrac{\Delta}{4a^{2}}>0$ donc $P(x)$ est  du signe de $a$.
\item Si $\Delta=0$ alors $P(x)=a\left(x+\dfrac{b}{2a}\right)^{2}$ qui est du signe de $a$ sauf pour $-\dfrac{b}{2a}$.
\item Si $\Delta>0$ alors $P(x)=a(x-r_{1})(x-r_{2})$ où $r_{1}$ et $r_{2}$ sont les racines de $P$. On peut donc dresser le tableau de signe suivant :

\end{itemize}
\vspace{1ex}
\begin{tabular}[c]{|c|ccccccc|} \hline
$x$ & $-\infty$ &  & $x_{1}$ &  & $x_{2}$ &  & $+\infty$ \\ \hline
$a$ &  & Signe de $a$ &  & Signe de $a$ &  & Signe de $a$ & \\ \hline
$x-x_{1}$ &  & $-$ & 0 & $+$ &  & $+$ & \\ \hline
$x-x_{2}$ &  & $-$ &  & $-$ & 0 & $+$ & \\ \hline
Signe de $P(x)$ &  & Signe de $a$ & 0 & Signe de $-a$ & 0 & Signe de $a$ & \\ \hline
\end{tabular}

\vspace{1ex}
\end{demo}

\begin{exemple}
Déterminer le tableau de signe de $P$ et $Q$ définis par
 \[P(x)=2x^{2}+4x+6 \quad \text{et} \quad Q(x)=-x^{2}-2x+1\]
\end{exemple}


Pour $P$ : $\Delta>0$ donc

\begin{center}
\begin{tabular}[c]{|c|ccc|} \hline
$x$ & $-\infty$ &  & $+\infty$ \\ \hline
Signe de $P(x)$ &  & + & \\ \hline
\end{tabular}
\end{center}

Pour $Q$ : Les racines sont $-1-\sqrt{2}$ et $-1+\sqrt{2}$ et de plus $-1<0$ donc

\vspace{1ex}
\begin{center}
\begin{tabular}[c]{|c|ccccccc|} \hline
$x$ & $-\infty$ &  & $-1-\sqrt{2}$ &  & $-1+\sqrt{2}$ &  & $+\infty$ \\ \hline
Signe de $P(x)$ &  & $-$ & 0 & $+$ & 0 & $-$ & \\ \hline
\end{tabular}
\end{center}

%
\subsection{Représentation graphique}
\begin{prop}
Soit $\oij$ un repère du plan.

La représentation graphique de $P$ est l'image par la translation de vecteur $\vect{u}\cp{-\frac{b}{2a}}{-\frac{\Delta}{4a}}$ de la parabole représentant la fonction $x\mapsto ax^{2}$.
L'axe de sym\'etrie d'une parabole repr\'esentative d'une fonction polyn\^ome du second degr\'e  du type $P(x)=ax^2+bx+c$ est $x=\dfrac{-b}{2a}$.
C'est donc une parabole de sommet $S\left(-\frac{b}{2a};P(-\frac{b}{2a})\right)$.
\end{prop}
Illustration :
\begin{center}
 \begin{tabular}{|c|c|c|c|}
   \hline
    & $\Delta >0$ & $\Delta =0$ & $\Delta <0$  \\
   \hline
   Factorisation de $P(x)$ & $a(x-x_{1})(x-x_{2})$ & $a(x-x_{0})^{2}$ &
   pas de factorisation  \\
   \hline
   Équation $P(x)=0$ & $2$ solutions $x_{1}$ et $x_{2}$ & une solution $x_{0}$ & pas de solution\\ \hline
   \rule[0pt]{0pt}{2.2cm} \raisebox{1.1cm}[0pt][0pt]{$a>0$} & \includegraphics[scale=0.7]{CoursSecondDegFigures.1} & \includegraphics[scale=0.7]{CoursSecondDegFigures.2}
   & \includegraphics[scale=0.7]{CoursSecondDegFigures.3} \\
   \hline
   \rule[0pt]{0pt}{2.2cm} \raisebox{1.1cm}[0pt][0pt]{$a<0$} & \includegraphics[scale=0.7]{CoursSecondDegFigures.4} & \includegraphics[scale=0.7]{CoursSecondDegFigures.5} &  \includegraphics[scale=0.7]{CoursSecondDegFigures.6} \\
   \hline
  \end{tabular}
  \end{center}

\end{document}